% =========================================================
% Proof Stub: Tail Suprema Are Monotone
% Source: volume-iii/analysis/sequences/notes/sequences/notes-k-tails.tex
% Mode: standard
% =========================================================

\clearpage
\phantomsection
\hypertarget{proof-RA-ABB-T-tail-sup-monotone}{}

\subsubsection[Tail Suprema Are Monotone]{Proof --- Tail Suprema Are Monotone}

\bigskip

\noindent
\textbf{Source.}
See \texttt{notes-k-tails.tex}.

\vspace{0.75em}

\noindent
\textbf{Goal.}
If $(x_n)$ is bounded and $s_k = \sup\{x_n : n \ge k\}$, then $(s_k)$ is decreasing.

\vspace{0.75em}

\noindent
\textbf{Logical form.}
\[
  T_k \supseteq T_{k+1} \Rightarrow s_k = \sup T_k \ge \sup T_{k+1} = s_{k+1}.
\]

\vspace{0.75em}

\noindent
\textbf{Proof strategy.}
Since $T_{k+1} \subseteq T_k$, every upper bound of $T_k$ is an upper bound of $T_{k+1}$, so the infimum of upper bounds (i.e.\ the supremum) can only decrease.

\vspace{0.75em}

\noindent
\textbf{Proof.}
\begin{proof}
\Claim If $(x_n)$ is bounded and $s_k = \sup\{x_n : n \ge k\}$, then $(s_k)$ is decreasing.

\Given 

\Goal 

\Strategy Since $T_{k+1} \subseteq T_k$, every upper bound of $T_k$ is an upper bound of $T_{k+1}$, so the infimum of upper bounds (i.e.\ the supremum) can only decrease.

\medskip

% --- observe T_{k+1} = {x_n : n ≥ k+1} ⊆ {x_n : n ≥ k} = T_k ---
% --- apply Monotonicity of Supremum: A ⊆ B ⇒ sup A ≤ sup B ---
% --- therefore s_{k+1} ≤ s_k, i.e. (s_k) is decreasing ---

\end{proof}

\begin{remark*}[Proof shape]
Nested tail sets give decreasing tail suprema by monotonicity of $\sup$.
\end{remark*}

\begin{remark*}[Dependencies]
\begin{itemize}
  \item Monotonicity of Supremum and Infimum (Bounding chapter, Theorem~\ref{thm:mono-sup-inf}).
  \item $(x_n)$ bounded ensures each $s_k$ exists as a finite real number.
\end{itemize}
\end{remark*}
